%----------------------------------------------------------------------------------------
%	科研版 · 中文 — 研究兴趣、学术论文、科研与工程经历、教学经历
%	叙事定位：研究问题、方法论、文献对比、学术服务
%----------------------------------------------------------------------------------------

%----------------------------------------------------------------------------------------
%	RESEARCH INTERESTS
%----------------------------------------------------------------------------------------

\section{研究兴趣}
高性能计算与 GPU 体系结构、算法设计与分析。

%----------------------------------------------------------------------------------------
%	PUBLICATIONS
%----------------------------------------------------------------------------------------

\section{学术论文}
\begin{tabularx}{\linewidth}{@{}X@{}}
\textbf{Disproofs of the Four-Coin Hereditary Subcurrency Conjecture and the Prefix Pattern Representative Conjecture.} \normalsize（在投 SODA 2027）研究换零钱问题中贪心算法保持最优的 orderly 货币系统，构造反例推翻了关于遗传性子货币模式的两项已有猜想。 \\
\end{tabularx}

%----------------------------------------------------------------------------------------
%	RESEARCH & ENGINEERING EXPERIENCE
%----------------------------------------------------------------------------------------

\section{科研与工程经历}

\begin{joblong}{阿里巴巴 | 推理引擎 · 算子优化实习生}{\DateAlibabaCn}
\item 参与 \textbf{PAI-vLLM}（阿里内部 vLLM 推理引擎）的算子优化，负责 \textbf{Qwen3-Next} 线性注意力（Gated DeltaNet）decode / MTP-verify 递推算子（Triton 实现）。
\item 定位到递推内核的瓶颈在于“输出串行依赖更新后的状态”，对递推做代数重写（$o_t = d\,(h_{t-1}^\top q_t) + u_t\,(k_t^\top q_t)$），使输出直接由更新前的状态算出：两次状态扫描合并为一次，输出计算脱离每步关键路径。
\item 设计 bf16 混合精度流水线（仅输入输出保留 fp32），并论证其数值安全性——decay 门控对累积误差构成压缩映射（~6e-3 且不随序列长度增长）；状态寄存器减半、SM occupancy 翻倍。
\item 建立评估方法（对照 fp64 参考实现的正确性验证、CUPTI 内核计时、K≠V 与非 2 幂形状），并借此发现 Triton autotune 基准测试污染内核状态的 bug；MTP-verify 内核提速 \textbf{10.9\%}，28/28 测试通过。
\end{joblong}

\begin{joblong}{摩尔线程 | GPU 计算与 AI 软件实习生}{\DateMooreCn}
\item 参与自研高性能 GPU 数学库 muBLAS（对标 NVIDIA cuBLAS）的核心算子研发，独立负责批处理稠密线性代数例程：LU 分解 \texttt{getrfBatched}、线性方程组求解 \texttt{getrsBatched}、矩阵求逆 \texttt{getriBatched} / \texttt{matinvBatched} 的设计与实现。
\item 系统分析共享内存布局、访存合并与归约策略对性能与数值稳定性的影响，通过 profiling 定位瓶颈并迭代优化内核（“建模—实验—验证”的调优方法论）；调研 cuBLAS / MAGMA 与相关文献，对比批处理小矩阵下选主元（pivoting）策略的精度—性能权衡。
\end{joblong}

\begin{joblong}{LemonDB：现代 C++ 多线程内存数据库}{\href{https://github.com/zty2004/ECE4821Jp2}{项目链接}}
\item 基于 \textbf{C++20} 实现高并发关系型数据库，设计拉取式线程池（批量拉取任务将全局锁竞争降低 \textbf{16 倍}）与依赖感知调度系统，自动解析查询间资源依赖并进行拓扑调度。
\end{joblong}

\begin{joblong}{MSD 分布式音乐推荐系统}{\href{https://github.com/zty2004/ECE4721J}{项目链接}}
\item 基于百万歌曲数据集构建端到端大数据管道（288GB 原始数据转为 884MB Avro，PySpark 处理耗时从 5.5 小时缩短至 \textbf{1.3 小时}），并用 FAISS 实现亚秒级近似最近邻推荐。
\end{joblong}

%----------------------------------------------------------------------------------------
%	TEACHING
%----------------------------------------------------------------------------------------

\section{教学经历}

\begin{jobshort}{CS4770J 算法导论 | 助教}{\DateTACn}
开展 OCaml 算法实验课（尾递归优化、延续传递风格）；维护并扩展课程 OJ 与评分脚本。
\end{jobshort}

\begin{jobshort}{CS2810J 数据结构与算法 | 助教}{\DateTACn}
主持习题课，讲授高级数据结构与复杂度分析；设计竞赛风格补充习题集并维护课程 OJ。
\end{jobshort}
